% TCC — Conclusão
\section{Conclusão}\label{sec:conclusoes}

Este trabalho investigou quais contribuições e limitações a combinação de
\textit{fuzzing} \textit{black-box} orientado por OpenAPI, deduplicação e
priorização \textit{fuzzy} apresenta para a triagem de divergências em APIs
REST. Para isso, foi implementado um fluxo que interpreta um subconjunto de
OpenAPI, gera casos-base e mutações, executa as requisições, identifica
divergências, agrupa observações equivalentes, atribui prioridade técnica e
preserva os artefatos de cada campanha.

Nas três dimensões avaliadas, os resultados indicam:

\begin{itemize}
    \item \textbf{execução e completude:} as 25 campanhas da coleta final
    concluíram as 7.175 requisições planejadas, e as contagens do artigo podem
    ser rastreadas até os planos, as respostas e os achados publicados;
    \item \textbf{estabilidade:} 441 dos 570 IDs semânticos reapareceram nas
    cinco repetições da respectiva API e 453 em pelo menos quatro. A presença
    foi completa no Petclinic REST e no Meilisearch; em VAmPI e RealWorld, a
    variação veio de dados gerados em cada campanha, e a crAPI concentrou a
    maior parte dos IDs observados uma única vez;
    \item \textbf{deduplicação e prioridade:} o agrupamento reduziu 4.000
    achados brutos a 2.460 assinaturas estritas (38,5\%). Na Etapa 1, o FIS
    revisado concordou com a referência cega em 120 de 122 famílias do
    Meilisearch, contra 4 da linha de base baseada somente no status.
\end{itemize}

A principal contribuição é um processo rastreável que transforma um contrato
OpenAPI em evidências organizadas para triagem de \textit{contract drift},
incompatibilidades funcionais, problemas de robustez e indícios de controle de
acesso, com uma prioridade cujos critérios podem ser auditados. Os resultados
não confirmam vulnerabilidades nem estimam impacto real. As conclusões são
limitadas pelo número de aplicações, pelos recortes dos contratos, pelo OpenAPI
como oráculo parcial e por uma referência de prioridade produzida por um modelo
de linguagem, sem casos \texttt{HIGH} no Meilisearch. Na coleta final, os
escores têm interpretação descritiva.

Como trabalhos futuros, ficam a ampliação das aplicações e dos recursos de
OpenAPI suportados, a resolução de dependências entre operações, a avaliação da
triagem por especialistas humanos, a integração com uma esteira de integração
contínua e a comparação com outras ferramentas sob um mesmo protocolo.
